\chapter*{APPENDIX}
\addcontentsline{toc}{chapter}{APPENDIX}

\section*{Appendix A: Staff Presence - ESP32 Firmware (C++)}
The full source code for the ESP32 node handled in the Arduino IDE. This handles dual-mode presence detection (active tag scanning and passive smartphone validation).

\begin{lstlisting}[language=C++]
/*
 * Staff Presence - ESP32 Firmware
 * V4: FULL DUAL-MODE (Hardware Tag Scanner + Mobile Verification Mailbox)
 *
 * This version supports BOTH ways for a staff member to be detected:
 * 1. AUTOMATIC: Scans for physical staff iBeacon tags (Averaged RSSI).
 * 2. MOBILE: Accepts a physical BLE connection from the staff phone.
 */

#include <WiFi.h>
#include <HTTPClient.h>
#include <BLEDevice.h>
#include <BLEUtils.h>
#include <BLEScan.h>
#include <BLEAdvertisedDevice.h>
#include <BLEServer.h>
#include <BLE2902.h>
#include <map>
#include <vector>
#include <numeric>

// Configuration
const char* ssid        = "Thorfinn";
const char* password    = "12345677";
const char* serverUrl   = "http://10.31.158.78:5000/api/ble-data";
const char* classroomId = "COMPUTERLAB";

#define SERVICE_UUID        "4fafc201-1fb5-459e-8fcc-c5c9c331914b"
#define CHARACTERISTIC_UUID "beb5483e-36e1-4688-b7f5-ea07361b26a8"

const int scanTime = 5; 
const int MIN_READINGS_TO_SEND = 1;
const int RSSI_THRESHOLD = -120; 

BLEScan* pBLEScan;
BLEServer* pServer;
bool mobileConnected = false;
int lastMobileRssi = -30;               
esp_bd_addr_t current_remote_bda = {0}; 
std::map<std::string, std::vector<int>> tagAccumulator;

void my_gap_event_handler(esp_gap_ble_cb_event_t event, esp_ble_gap_cb_param_t *param) {
    if (event == ESP_GAP_BLE_READ_RSSI_COMPLETE_EVT) {
        if (param->read_rssi_cmpl.status == ESP_BT_STATUS_SUCCESS) {
            lastMobileRssi = param->read_rssi_cmpl.rssi;
            Serial.printf("[Signal] Live Mobile RSSI: %d dBm\n", lastMobileRssi);
        }
    }
}

void reportToBackend(String staffUuid, int rssi, String method) {
    if (WiFi.status() != WL_CONNECTED) {
        Serial.printf("[Offline Mode] Saw %s but WiFi is down.\n", staffUuid.c_str());
        return;
    }
    HTTPClient http;
    http.begin(serverUrl);
    http.addHeader("Content-Type", "application/json");

    String payload = "{\"esp32_id\":\"" + String(classroomId) +
                     "\",\"beacon_uuid\":\"" + staffUuid +
                     "\",\"rssi\":" + String(rssi) + 
                     ",\"method\":\"" + method + "\"}";

    int code = http.POST(payload);
    Serial.printf("[%s] %s | Avg RSSI: %d | HTTP: %d\n", method.c_str(), staffUuid.c_str(), rssi, code);
    http.end();
}

class MobileMailbox : public BLECharacteristicCallbacks {
    void onWrite(BLECharacteristic *pChar) {
        String val = pChar->getValue();
        if (val.length() > 0) {
            String staffUuid = val;
            staffUuid.toUpperCase();
            
            esp_ble_gap_read_rssi(current_remote_bda);
            
            Serial.printf("[Mobile Verification] Staff: %s | Real RSSI: %d\n", staffUuid.c_str(), lastMobileRssi);
            reportToBackend(staffUuid, lastMobileRssi, "mobile_verification");
        }
    }
};

class ServerCallbacks : public BLEServerCallbacks {
    void onConnect(BLEServer* pServer, esp_ble_gatts_cb_param_t *param) {
        memcpy(current_remote_bda, param->connect.remote_bda, 6);
        mobileConnected = true;
        Serial.println("Mobile Connected.");
        esp_ble_gap_read_rssi(current_remote_bda); 
    }
    
    void onDisconnect(BLEServer* pServer) {
        mobileConnected = false;
        Serial.println("Mobile Disconnected.");
        pServer->getAdvertising()->start();
    }
};

class TagScanCallback : public BLEAdvertisedDeviceCallbacks {
    void onResult(BLEAdvertisedDevice dev) {
        int rssi = dev.getRSSI();
        String devName = dev.getName().c_str();
        String devAddr = dev.getAddress().toString().c_str();

        if (rssi < RSSI_THRESHOLD) return;

        if (devName.length() > 0) {
            Serial.printf("[Scanner] Found: \"%s\" (Addr: %s) | RSSI: %d\n", devName.c_str(), devAddr.c_str(), rssi);
        }

        uint8_t* pay = dev.getPayload();
        size_t len = dev.getPayloadLength();

        for (int i = 0; i < (int)len - 20; i++) {
            if (pay[i] == 0x4C && pay[i+1] == 0x00 && pay[i+2] == 0x02 && pay[i+3] == 0x15) {
                char buf[33];
                for (int j = 0; j < 16; j++) sprintf(&buf[j*2], "%02X", pay[i+4+j]);
                tagAccumulator[std::string(buf)].push_back(rssi);
                
                Serial.printf("iBeacon Found! UUID: %s | RSSI: %d\n", buf, rssi);
                return;
            }
        }
    }
};

void setup() {
    Serial.begin(115200);

    WiFi.begin(ssid, password);
    while (WiFi.status() != WL_CONNECTED) { delay(500); Serial.print("."); }
    Serial.println("\nWiFi OK.");

    BLEDevice::init(classroomId);
    BLEDevice::setCustomGapHandler(my_gap_event_handler);

    pServer = BLEDevice::createServer();
    pServer->setCallbacks(new ServerCallbacks());
    BLEService *pService = pServer->createService(SERVICE_UUID);
    BLECharacteristic *pChar = pService->createCharacteristic(CHARACTERISTIC_UUID, BLECharacteristic::PROPERTY_WRITE);
    pChar->setCallbacks(new MobileMailbox());
    pService->start();

    BLEAdvertising *pAdv = BLEDevice::getAdvertising();
    pAdv->addServiceUUID(SERVICE_UUID);
    pAdv->setScanResponse(true);
    pAdv->start();

    pBLEScan = BLEDevice::getScan();
    pBLEScan->setAdvertisedDeviceCallbacks(new TagScanCallback(), false);
    pBLEScan->setActiveScan(true);
    pBLEScan->setInterval(100);
    pBLEScan->setWindow(99);

    Serial.println("Dual-Mode Active: Scanning for Tags + Waiting for Mobile.");
}

void loop() {
    Serial.println("\n--- BLE Scan Window (Scanning for Tags) ---");
    tagAccumulator.clear();
    
    BLEScanResults* results = pBLEScan->start(scanTime, false);
    pBLEScan->clearResults();

    for (auto& entry : tagAccumulator) {
        const std::string& uuid = entry.first;
        std::vector<int>&  readings = entry.second;

        if ((int)readings.size() >= MIN_READINGS_TO_SEND) {
            int sum = std::accumulate(readings.begin(), readings.end(), 0);
            int avgRssi = sum / (int)readings.size();
            reportToBackend(String(uuid.c_str()), avgRssi, "hardware_tag");
        }
    }

    delay(2000); 
}
\end{lstlisting}

\newpage
\section*{Appendix B: Context-Aware Business Logic (Node.js)}
The core logic for parsing the timetable, checking substitutions, and calculating proximity.

\begin{lstlisting}[language=JavaScript]
const Staff = require('../models/Staff');
const Timetable = require('../models/Timetable');
const Notification = require('../models/Notification');
const Attendance = require('../models/Attendance');
const User = require('../models/User');
const Leave = require('../models/Leave');
const { sendEmail } = require('../utils/emailService');

const getDayName = () => {
    const days = ['Sunday', 'Monday', 'Tuesday', 'Wednesday', 'Thursday', 'Friday', 'Saturday'];
    return days[new Date().getDay()];
};

const getCurrentTime = () => {
    const now = new Date();
    return `${now.getHours().toString().padStart(2, '0')}:${now.getMinutes().toString().padStart(2, '0')}`;
};

async function isStaffPermitted(staffId, classroomId) {
    const now = new Date();
    const days = ['Sunday', 'Monday', 'Tuesday', 'Wednesday', 'Thursday', 'Friday', 'Saturday'];
    const currentDay = days[now.getDay()];
    const currentTime = now.toTimeString().split(' ')[0].substring(0, 5);
    const startOfToday = new Date();
    startOfToday.setHours(0, 0, 0, 0);

    // 1. Direct match in timetable
    const directSlot = await Timetable.findOne({
        staff_id: staffId,
        classroom_id: classroomId,
        day_of_week: currentDay,
        start_time: { $lte: currentTime },
        end_time: { $gte: currentTime }
    });
    if (directSlot) return { permitted: true, slot: directSlot, type: 'direct' };

    // 2. Check for substitution (Was someone else supposed to be here?)
    const scheduledSlot = await Timetable.findOne({
        classroom_id: classroomId,
        day_of_week: currentDay,
        start_time: { $lte: currentTime },
        end_time: { $gte: currentTime }
    });

    if (scheduledSlot) {
        const scheduledStaffId = scheduledSlot.staff_id;
        
        const onLeave = await Leave.findOne({
            staff_id: scheduledStaffId,
            status: 'approved',
            start_date: { $lte: now },
            end_date: { $gte: now }
        });

        const presentToday = await Attendance.findOne({
            staff_id: scheduledStaffId,
            date: { $gte: startOfToday }
        });

        if (onLeave || !presentToday) {
            const user = await User.findOne({ staff_id: staffId });
            if (user) {
                const subNotification = await Notification.findOne({
                    recipient_id: user._id,
                    type: { $in: ['meeting_request', 'substitution_request'] },
                    'related_data.classroomId': classroomId.toString()
                }).sort({ createdAt: -1 });

                if (subNotification) return { permitted: true, slot: scheduledSlot, type: 'substitution' };
            }
        }
    }
    return { permitted: false };
}

exports.getAllStaffStatus = async (req, res) => {
    try {
        const staffMembers = await Staff.find().lean();
        const currentDay = getDayName();
        const currentTime = getCurrentTime();
        
        const startOfToday = new Date();
        startOfToday.setHours(0, 0, 0, 0);
        const endOfToday = new Date();
        endOfToday.setHours(23, 59, 59, 999);

        const [allAttendanceToday, allActiveClasses, allApprovedLeaves, allLatestAttendance] = await Promise.all([
            Attendance.find({ date: { $gte: startOfToday, $lte: endOfToday } })
                .populate('classroom_id', 'room_name').lean(),
            Timetable.find({ 
                day_of_week: currentDay, 
                start_time: { $lte: currentTime }, 
                end_time: { $gte: currentTime } 
            }).populate('classroom_id', 'room_name').lean(),
            Leave.find({ 
                status: 'approved', 
                start_date: { $lte: endOfToday }, 
                end_date: { $gte: startOfToday } 
            }).lean(),
            Attendance.aggregate([
                { $sort: { check_in_time: -1 } },
                { $group: {
                    _id: '$staff_id',
                    latest_record: { $first: '$$ROOT' }
                }},
                { $lookup: {
                    from: 'classrooms',
                    localField: 'latest_record.classroom_id',
                    foreignField: '_id',
                    as: 'room'
                }},
                { $unwind: { path: '$room', preserveNullAndEmptyArrays: true } }
            ])
        ]);

        const attendanceTodayMap = new Map();
        allAttendanceToday.forEach(a => {
            const sid = (a.staff_id || '').toString();
            if (!sid) return;
            const existing = attendanceTodayMap.get(sid);
            if (!existing || new Date(a.last_seen_time || a.check_in_time) > new Date(existing.last_seen_time || existing.check_in_time)) {
                attendanceTodayMap.set(sid, a);
            }
        });

        const latestAttendanceMap = new Map();
        allLatestAttendance.forEach(a => {
            if (a._id) latestAttendanceMap.set(a._id.toString(), a);
        });

        const activeClassMap = new Map();
        allActiveClasses.forEach(c => activeClassMap.set((c.staff_id || '').toString(), c));

        const leaveMap = new Map();
        allApprovedLeaves.forEach(l => leaveMap.set((l.staff_id || '').toString(), l));

        const now = new Date();
        const signalThreshold = 5 * 60 * 1000; 

        const staffStatus = staffMembers.map(staff => {
            const sid = staff._id.toString();
            const attendanceToday = attendanceTodayMap.get(sid);
            const latestHistory = latestAttendanceMap.get(sid);
            const activeClass = activeClassMap.get(sid);
            const approvedLeave = leaveMap.get(sid);

            const isStale = attendanceToday && (now - new Date(attendanceToday.last_seen_time || attendanceToday.check_in_time) > signalThreshold);

            let liveStatus = 'Absent';
            if (approvedLeave) {
                liveStatus = 'On Leave';
            } else if (attendanceToday && !isStale) {
                liveStatus = attendanceToday.status; 
            } else if (attendanceToday && isStale) {
                liveStatus = 'Left'; 
            }

            const latestRec = latestHistory?.latest_record;

            return {
                ...staff,
                currentStatus: liveStatus,
                lastSeen: latestRec ? (latestRec.last_seen_time || latestRec.check_in_time) : null,
                lastSeenLocation: latestHistory?.room?.room_name || 'unknown',
                currentLocation: (attendanceToday && !isStale) ? attendanceToday.classroom_id.room_name : 
                                 (latestRec ? `Last seen in ${latestHistory.room?.room_name || 'unknown'}` : 'Not detected'),
                expectedLocation: activeClass ? activeClass.classroom_id.room_name : 'No Class Assigned',
                isCorrectLocation: activeClass && attendanceToday && !isStale ?
                    (activeClass.classroom_id?._id?.toString() === attendanceToday.classroom_id?._id?.toString()) :
                    (activeClass ? false : true),
                activeClass: activeClass ? {
                    subject: activeClass.subject,
                    room: activeClass.classroom_id?.room_name,
                    startTime: activeClass.start_time,
                    endTime: activeClass.end_time
                } : null
            };
        });

        res.status(200).json(staffStatus);
    } catch (error) {
        console.error('Error in getAllStaffStatus:', error);
        res.status(500).json({ message: 'Server error fetching staff status.' });
    }
};
\end{lstlisting}

\newpage
\section*{Appendix C: Main Server API Routes (Node.js)}
Initialization of the MERN stack application, caching setup, and endpoint configurations.

\begin{lstlisting}[language=JavaScript]
require('dotenv').config();
const express = require('express');
const bodyParser = require('body-parser');
const cors = require('cors');
const morgan = require('morgan');
const cron = require('node-cron');
const compression = require('compression');
const multer = require('multer');
const path = require('path');
const fs = require('fs');
const connectDB = require('./config/database');

const Staff = require('./models/Staff');
const Classroom = require('./models/Classroom');
const Timetable = require('./models/Timetable');
const Attendance = require('./models/Attendance');
const Alert = require('./models/Alert');
const User = require('./models/User');

const { generateToken, authenticateToken, requireAdmin } = require('./middleware/auth');
const app = express();
const PORT = process.env.PORT || 5000;

connectDB();

app.use(cors());
app.use(compression());
app.use(morgan('dev'));
app.use(bodyParser.json());

const staffCache = new Map(); 
const classroomCache = new Map(); 
const bleThrottle = new Map(); 
const rssiBuffer = new Map(); 
const RSSI_BUFFER_SIZE = 5;   

setInterval(() => {
    staffCache.clear();
    classroomCache.clear();
    console.log('⚡ Performance: Cleared Staff/Classroom ID caches.');
}, 10 * 60 * 1000);

setInterval(() => {
    bleThrottle.clear();
    rssiBuffer.clear();
    console.log('⚡ Performance: Reset BLE throttle & RSSI buffer maps.');
}, 60 * 60 * 1000);

app.post('/api/auth/login', async (req, res) => {
    try {
        let { email, password } = req.body;

        if (email) {
            email = email.trim().toLowerCase();
        }

        console.log(`Login attempt for: ${email}`);
        const user = await User.findOne({ email }).populate('staff_id');

        if (!user) {
            return res.status(401).json({ error: 'Invalid credentials' });
        }

        if (!user.isActive) {
            return res.status(401).json({ error: 'Account is deactivated' });
        }

        const isMatch = await user.comparePassword(password);

        if (!isMatch) {
            return res.status(401).json({ error: 'Invalid credentials' });
        }

        const token = generateToken(user._id, user.role);

        res.json({
            token,
            user: {
                id: user._id,
                email: user.email,
                name: user.name,
                role: user.role,
                staff_id: user.staff_id
            }
        });
    } catch (err) {
        res.status(500).json({ error: err.message });
    }
});

app.listen(PORT, () => console.log(`[Express] Server running on port ${PORT}`));
\end{lstlisting}
